%% ma: L10-B01-0037
%% lop: 10
%% chuong: 1
%% bai: 1
%% dang: 10.1.2
%% loai: TLN
%% muc: TH
%% dap_an: "2"
%% nguon: "Ảnh giáo viên gửi 10/10/2026 (ThS. Nguyễn Hoàng Việt, Chương 1 Mệnh đề), B-Đề 2, Phần 3 Câu 21"
%% kien_thuc: "Bài 1 (mệnh đề, mệnh đề chứa biến)"
%% trang_thai: cho-duyet
%% ly_do: "Dạng toán mới đề xuất 10.1.2: xét tính đúng sai của mệnh đề (số học, đại số, hình học)"
%% ghi_chu: "Đề gốc không có đáp án. Mệnh đề cuối thuộc kiến thức địa lí (không phải Toán), giáo viên quyết định có giữ không"
\begin{ex}
Cho các mệnh đề sau. Hãy cho biết có bao nhiêu mệnh đề sai.
\begin{itemize}
\item Phương trình $x^2-3x+8=0$ có nghiệm.
\item $16$ không là số nguyên tố.
\item Hai phương trình $x^2-4x+3=0$ và $x^2-\sqrt{x+3}+1=0$ có nghiệm chung.
\item Buôn Ma Thuột là thành phố của tỉnh Quảng Ngãi.
\end{itemize}
\shortans{2}
\loigiai{Mệnh đề 1 sai ($\Delta=9-32<0$). Mệnh đề 2 đúng ($16=2^4$). Mệnh đề 3 đúng: $x^2-4x+3=0$ có nghiệm $x=1$ hoặc $x=3$; $x=1$ thỏa $1-\sqrt4+1=0$. Mệnh đề 4 sai (Buôn Ma Thuột thuộc tỉnh Đắk Lắk). Có $2$ mệnh đề sai.}
\end{ex}
